\begin{helper}[Terminal parallel preprocessing]
Let $T$ be a positive-rung maximum ETT whose terminal step is parallel,
and assume the induction context from condition (5.1): lower-rung
elementarity and interchangeability, stability descent and disjoint stored
pairs, the terminal-component intersection assertion, the defective-endpoint
order assertion, and the mod-coloring assertion, together with current-rung
assertions (iii)--(vi).  Let $T_n$ be its terminal
ladder tree, let $\sigma_n$ be a proper
$(T_n,D_n,\varphi_n)$-stable coloring leaving exactly $e_0$ uncolored, and
let $J_n$ be the support-prefix closure in a parallel hierarchy base.
Then there are the terminal support $v_n$, a vertex $w\ne v_n$ in both
$T_n$ and $J_n$, and terminal colors
$\gamma_n,\delta_n$ with their exact selected-field identities such that
$v_n\in V(T_n)\cap V(J_n)$, $\delta_n$ is missing at $v_n$ under
$\sigma_n$, $\gamma_n$ is missing nowhere on $T_n$ under $\sigma_n$,
$w$ has at least one missing color, and $V(T_n)$ is elementary under
$\sigma_n$.
\end{helper}
\begin{proof}
Put $n=r(T)>0$, and write $T_n$ for the terminal ladder tree.  By
\noderef{TashkinovSeries}, the last positive slot has a certified step.
The displayed PE-kind equality rules out the other constructors of
\noderef{TashkinovSeriesStep}.  The remaining constructor supplies a
maximum-defective choice
\[
 (\pi_{n-1},v_n,f_n,u_n,\delta_n),
\]
a second coloring $\rho=\pi'_{n-1}$, and a color $\gamma_n$.  Its exact
chosen-vertex, gamma, and delta equations are the three selected-field
identities required in the conclusion, and its coloring equation is
\[
 \varphi_n=\rho/P_{v_n}(\gamma_n,\delta_n,\rho).                 \tag{1}
\]
The fields of \noderef{MaximumDefectiveChoice}, together with
\noderef{DefectiveVertex}, put $v_n$ in $V(T_n)$ and say that $f_n$ is an
exact defective boundary-edge identity incident with $v_n$, has outside
endpoint $u_n$, and satisfies $\pi_{n-1}(f_n)=\delta_n$.  The rule
\noderef{ParallelExtensionColor} gives
$\gamma_n\in\overline\pi_{n-1}(v_n)$.  The second-stable relation in the
PE certificate preserves this missing set at $v_n$.  Also, $f_n$ is
incident with $T_n$ by \noderef{TreeIncidentEdge}, and $\delta_n$ is in
the enlarged control set, so the edge clause of
\noderef{StableColoring} gives
\[
 \rho(f_n)=\pi_{n-1}(f_n)=\delta_n.                             \tag{2}
\]
Thus \noderef{PresentColors} makes $\delta_n$ present at $v_n$ under
$\rho$, whereas the complementary definition \noderef{MissingColors}
makes $\gamma_n$ absent there.  In particular,
$\gamma_n\ne\delta_n$.

We next prove pre-switch elementarity.  If $n>1$, apply
\noderef{CertifiedETTRestriction} with $r=n-1$.  The restriction has
positive rung $r<n$, displayed tree exactly $T_n$, generating coloring
$\varphi_{n-1}$, and maximum property.  Instantiating the
\texttt{lowerOne} field of \noderef{TechnicalInductionContext} with that
actual restricted ETT makes $V(T_n)$ elementary under
$\varphi_{n-1}$.  If $n=1$, the initial certificate retained by the
series and \noderef{TechnicalTheoremBase} give the same assertion.  The
selected-coloring field makes $\pi_{n-1}$ stable relative to
$\varphi_{n-1}$, and the PE second-stable field makes $\rho$ stable
relative to $\pi_{n-1}$.  The vertex clause of
\noderef{StableColoring} therefore identifies the three missing-color
sets at every vertex of $T_n$.  Expanding
\noderef{ColoredVertexSetElementary}, $V(T_n)$ is elementary under both
$\pi_{n-1}$ and $\rho$.

The PE constructor says that $V(T_n)$ is closed under
$\varphi_{n-1}$.  It is also closed under $\rho$.  Indeed, suppose that
an exact identity $a$ were a boundary edge under
\noderef{BoundaryEdges}, that $\rho(a)=c$, and that
$c\in\overline\rho(T_n)$.  A boundary edge is incident with $T_n$ by
\noderef{TreeIncidentEdge}.  The vertex equalities just obtained put $c$
in the tree-missing set for each of the preceding colorings.  All three
colorings are proper and leave exactly $e_0$ uncolored by the PE and
maximum-choice fields.  Symmetry and transitivity in \noderef{sc2},
applied first to $\rho$ and $\pi_{n-1}$ and then to
$\pi_{n-1}$ and $\varphi_{n-1}$, fix the controlled incident identity
$a$ successively.  Hence $\varphi_{n-1}(a)=c$, contradicting
\noderef{ColoredVertexSetClosed}.  Equation (2) and the fact that $f_n$
is a boundary identity now show
\[
 \delta_n\notin\overline\rho(T_n).                              \tag{3}
\]

We spell out the application of \noderef{ExtensionBase}.  Given any
positive-rung maximum ETT $\mathcal E'$ on the fixed graph, palette, and
edge $e_0$, with arbitrary input coloring and rung $0<s<n$, instantiate
the six universally quantified fields of
\noderef{TechnicalInductionContext} on $\mathcal E'$.  The fields
\texttt{lowerOne} and \texttt{lowerTwo} are exactly its elementarity and
interchangeability premises; the two clauses of \texttt{lowerThree} are
its stability-descent and disjoint-pair premises; \texttt{lowerFour} is
its terminal-component premise; \texttt{lowerFive} is its ordered
defective-endpoint premise; and \texttt{lowerSix} is its mod-coloring
premise.  Positivity, the strict bound, and maximum property are precisely
the common hypotheses of those six fields, so this supplies the entire
arbitrary-input lower package required by \noderef{ExtensionBase}.
For the present ETT, use the assumed positive rung, maximum property, and
terminal PE equation.  The PE constructor makes $\varphi_n$ proper and
exactly-$e_0$-uncolored, and reflexivity in \noderef{sc2} makes it
$(T_n,D_n,\varphi_n)$-stable.  Thus all remaining premises of
\noderef{ExtensionBase} are discharged, and its conclusion says that the
$(\gamma_n,\delta_n)$-component through $v_n$ under $\varphi_n$ meets
$V(T_n)$ exactly at $v_n$.

Equation (1) swaps the whole component through $v_n$.  Unfolding
\noderef{TwoColorComponentSwap} and
\noderef{TwoColorIncidenceGraph}, this changes only the labels
$\gamma_n,\delta_n$ inside that component and preserves every adjacency
of the two-color incidence graph.  Hence the support defined by
\noderef{TwoColorComponent} is the same under $\rho$ and $\varphi_n$,
and its intersection with $V(T_n)$ is still the singleton $\{v_n\}$.
Moreover, (2), missingness of $\gamma_n$ at $v_n$, adjacency as defined
in \noderef{AdjacentEdges}, and properness in
\noderef{IsProperEdgeColoring} show that $f_n$ is the unique
$\gamma_n$- or $\delta_n$-colored edge incident with $v_n$ under $\rho$.
If an edge-node of this component were incident with $T_n$, choose its
tree endpoint using \noderef{TreeIncidentEdge}; the incidence rule puts
that endpoint in the same component, so it is $v_n$, and the edge is
$f_n$.  Consequently (1) fixes every tree-incident edge except $f_n$,
changes $f_n$ from $\delta_n$ to $\gamma_n$, and makes $\delta_n$
missing at $v_n$ under $\varphi_n$.

For vertices other than $v_n$, the missing sets under $\rho$ and
$\varphi_n$ are therefore equal.  At $v_n$, the latter is obtained from
the former by deleting $\gamma_n$ and inserting $\delta_n$.  The former
vertex sets are pairwise disjoint because $V(T_n)$ is elementary under
$\rho$, and (3) says that no vertex of $T_n$ missed $\delta_n$ under
$\rho$.  The definition \noderef{ColoredVertexSetElementary} now shows
that $V(T_n)$ is elementary under $\varphi_n$.  The assumed stability of
$\sigma_n$ identifies every vertex missing set with its
$\varphi_n$ counterpart.  It follows that $V(T_n)$ is elementary under
$\sigma_n$ and that
$\delta_n\in\overline\sigma_n(v_n)$.

It remains to find the common vertex.  The exact selected identity $f_n$
is defective.  By \noderef{DefectiveEdge} and
\noderef{DefectiveColor}, its color $\delta_n$ has a second, distinct
boundary identity $g$ under $\pi_{n-1}$.  Stability fixes both
$\delta_n$-colored identities in passing to $\rho$.  The identity $g$
cannot be in the switched component: its inside endpoint would then be
$v_n$, and $g$ and $f_n$ would be distinct adjacent edges of the same
color, contrary to \noderef{AdjacentEdges} and
\noderef{IsProperEdgeColoring}.  Thus (1) fixes $g$, so $g$ remains an
exact $\delta_n$-boundary identity under $\varphi_n$.  The delta-boundary
clause of \noderef{HkuParallelBoundaryStructure} supplies its inside
endpoint $w$ with $w\prec_{T_n}v_n$.  By
\noderef{TreeVertexOrder} and \noderef{TreeSequencePrefix},
$w\ne v_n$ and both $w$ and $v_n$ belong to the ordered prefix
$T_n(v_n)$.

The PE clause of \noderef{HierarchyBase} identifies that support position
with $v_n$ and makes $J_n$ a \noderef{TashkinovClosure} of this exact
prefix under $\sigma_n$.  Each step in its
\noderef{TashkinovAugmentationSequence} has the old tree as an ordered
initial segment by \noderef{TashkinovAugmentation}; transitivity of
\noderef{TreeSequenceInitialSegment}, together with
\noderef{TreeSequenceVertexSet}, therefore retains every prefix vertex.
Hence $v_n,w\in V(J_n)$, while their construction already puts them in
$V(T_n)$.  Inspecting the initial certificate of
\noderef{TashkinovSeriesStep} gives the original
\noderef{KTriple}, whose degree clause is $\Delta(G)+1\leq k$.
Applying \noderef{MissingColorsNonemptyOfMaximumDegree} at $w$ gives
$\overline\sigma_n(w)\ne\varnothing$.

Finally, \noderef{HkuParallelBoundaryStructure} gives the exact nonempty
identity equality
$\partial_{\varphi_n,\gamma_n}(T_n)=\{f_n\}$.  Its last clause, using
the already proved pre-step elementarity and
$\gamma_n\ne\delta_n$, would make this boundary empty if
$\gamma_n\in\overline\varphi_n(T_n)$, where tree missingness is the
union in \noderef{TreeMissingColors}.  Therefore $\gamma_n$ is not
tree-missing under $\varphi_n$.  The vertex clause of the assumed stable
relation transfers this to $\sigma_n$.  The vertices and colors just
constructed satisfy every assertion in the conclusion.
\end{proof}
